\subsection{深度与正则序列}

设 A 为环，M 为 A-模。回忆（A, X, 第 158 页），A 中元素的有限序列 \((x_{1},\ldots,x_{r})\) 称为关于 M 正则的，或 M-正则的，如果对 \(i=1,\ldots,r\) ，乘以 \(x_{i}\) 的映射在 A-模 \(\mathrm{M}/(x_{1}\mathrm{M}+\ldots+x_{i-1}\mathrm{M})\) 中是单射。设 \((x_{1},\ldots,x_{r})\) 为 M-正则序列；对每个平坦 A-模 N，序列 \((x_{1},\ldots,x_{r})\) 是 \(M\otimes_{A}N\) -正则的。若 \(\rho:A\to B\) 是使 B 成为平坦 A-模的环同态，则序列 \((\rho(x_{1}),\ldots,\rho(x_{r}))\) 对 B-模 \(B\otimes_{A}M\) 正则。特别地，对 A 的每个素理想 p，在 \(A_{p}\) 中序列 \((x_{1},\ldots,x_{r})\) 的像是 \(M_{p}\) -正则的。

在下文中我们主要考虑 M-正则序列这一概念在下列情形的形态：环 A 为局部 Noether 环，模 M 有限生成，且序列的元素属于 \(m_{A}\) ；此时 M-正则序列的概念与 M 的完全割序列的概念一致 (A, X, 第 160 页, 推论 1)。

命题 7.--- 设 A 为环，J 为 A 的理想，M 为 A-模，\((x_{1},\ldots,x_{r})\) 为由 J 中元素组成的 M-正则序列。则有

\[
\operatorname{prof} _ {\mathrm{A}} (\mathrm{J}; \mathrm{M}) = r + \operatorname{prof} _ {\mathrm{A}} (\mathrm{J}; \mathrm{M} / (x _ {1} \mathrm{M} + \dots + x _ {r} \mathrm{M}))
\]

特别地 \(\operatorname{prof}_{\mathrm{A}}(\mathrm{J};\mathrm{M})\geqslant r\)

\(r=1\) 的情形由第 1 节的注 5 应用于正合序列而得出

\[
0 \rightarrow \mathrm{M} \xrightarrow {(x _ {1}) _ {\mathrm{M}}} \mathrm{M} \longrightarrow \mathrm{M} / x _ {1} \mathrm{M} \rightarrow 0.
\]

一般情形由对 r 作归纳法从它推出。

定理 2.--- 设 A 为 Noether 环，J 为 A 的理想，M 为有限生成 A-模。

\begin{enumerate}
\def\labelenumi{\alph{enumi})}
\item
  设 \(\mathrm{prof}_{\mathrm{A}}(\mathrm{J};\mathrm{M})\) 有限。则每个由 J 中元素组成的 M-正则序列都可延拓为一个由 J 中元素组成、长度为 \(\mathrm{prof}_{\mathrm{A}}(\mathrm{J};\mathrm{M})\) 的 M-正则序列。
\item
  M 关于 J 的深度是由 J 中元素组成的 M-正则序列的长度的上确界。
\item
  \(\operatorname{prof}_{\mathrm{A}}(\mathrm{J};\mathrm{M})\) 有限必须且只需 M 的支撑集与 \(\mathrm{V}(\mathrm{J})\) 相交，或等价地，\(\mathrm{JM} \neq \mathrm{M}\) 。
\end{enumerate}

设 \((x_{1},\ldots,x_{r})\) 为由 J 中元素组成的 M-正则序列。有 \(r\leqslant\operatorname{prof}_{\mathrm{A}}(\mathrm{J};\mathrm{M})\) （命题 7）；设该不等式是严格的。用 N 表示 A-模 \(\mathrm{M}/(x_{1}\mathrm{M}+\ldots+x_{r}\mathrm{M})\) 。有 \(\operatorname{prof}_{\mathrm{A}}(\mathrm{J};\mathrm{N})>0\) （同上引文），故存在 J 的元素 x 使乘以 x 的映射 \(x_{N}\) 是单射（第 1 节, 注 2），即序列 \((x_{1},\ldots,x_{r},x)\) 是 M-正则的。由归纳法得出：对每个满足 \(r\leqslant s\leqslant\operatorname{prof}_{\mathrm{A}}(\mathrm{J};\mathrm{M})\) 的整数 s，序列 \((x_{1},\ldots,x_{r})\) 可补全为长度 s 的 M-正则序列，这就蕴含断言 a) 与 b)。断言 c) 由第 1 节的注 1 及第 3 节的定理 1 的推论 1 得出。

推论 1. 对每个由 J 中元素组成的 M-正则序列 \((x_{1},\ldots ,x_{r})\) ，下列性质等价：

\begin{enumerate}
\def\labelenumi{(\roman{enumi})}
\item
  有 \(r = \mathrm{prof}_{\mathrm{A}}(\mathrm{J};\mathrm{M})\)
\item
  序列 \((x_{1},\ldots,x_{r})\) 在由 J 中元素组成的 M-正则序列之中是极大的；
\item
  A-模 \(\mathrm{M}/(x_{1}\mathrm{M}+\ldots+x_{r}\mathrm{M})\) 有一个被 J 零化的非零元素；
\end{enumerate}

\[
\text {(iv)} \text {   we have   } \operatorname{Ass} \bigl (\mathrm{M} / (x _ {1} \mathrm{M} + \ldots + x _ {r} \mathrm{M}) \bigr) \cap \mathrm{V} (\mathrm{J}) \neq \emptyset .
\]

 与 的等价性由定理 2 得出；、 与 的等价性由第 1 节的注 2 应用于 A-模 M/(x₁M + \ldots{} + xᵣM) 而得出。

推论 2. --- 设 A 为 Noether 局部环，M 为非零的有限生成 A-模。则

\[
\operatorname{prof} _ {\mathrm{A}} (\mathrm{M}) \leqslant \dim_ {\mathrm{A}} (\mathrm{M}) <   + \infty .
\]

事实上，每个由 \(\mathfrak{m}_{\mathrm{A}}\) 中元素组成的 M-正则序列对 M 是完全割的 (A, X, 第 157 页, 命题 5)，从而对 M 是割的 (VIII, § 3, 第 2 节, 命题 3 的推论)；因此其长度以整数 \(\dim_{\Lambda}(\mathrm{M})\) 为上界（同上引文, 定理 1）。

